% \newpage
% \section{Brainstorming \& Problem Framing}

% \subsection{Background and the Paradigm Shift in Digital Content}
% The landscape of digital media and content creation has undergone a massive, disruptive paradigm shift with the rapid commercialization of Generative Artificial Intelligence (GenAI). Historically, producing high-fidelity visual, auditory, and textual content required significant human labor, specialized technical skills, and substantial time investment. Today, advanced multimodal GenAI models—such as Midjourney and Stable Diffusion for imagery, Google's Veo and OpenAI's Sora for video, and highly sophisticated voice cloning engines—can generate professional-grade, hyper-realistic content almost instantaneously based on simple natural language prompts. \\

% To thoroughly dissect the multidimensional implications of this technological leap, our team utilized structured brainstorming methods, specifically Mind Mapping and the 5W1H (Who, What, When, Where, Why, How) framework. Through this rigorous ideation phase, we recognized that while GenAI exponentially democratizes creative capabilities, its foundational engineering architecture is inherently controversial. These models rely on algorithms (such as latent diffusion and large language models) that require training on incredibly massive datasets. Consequently, tech companies have been indiscriminately harvesting, or ``scraping,'' billions of data points from the public internet. This practice has sparked an intense, unprecedented global debate regarding the ethical acquisition of training data, the boundaries of copyright law, and the profound societal consequences of indistinguishable synthetic media.

% \subsection{Core Problem Definition}
% The core engineering and ethical problem lies at the volatile intersection of rapid technological innovation, intellectual property (IP) appropriation, and the preservation of digital truth. Through our analysis, we distilled the problem into two primary dilemmas:
% \begin{itemize}
%     \item \textbf{Data Provenance and Copyright Infringement:} Current state-of-the-art AI models are built upon billions of copyrighted images, artworks, literary works, and videos scraped from the web. This is done without the explicit consent, credit, or financial compensation of the original creators. Tech companies often shield this practice behind the legal argument of ``Fair Use'' or text-and-data mining exceptions, arguing that algorithms are simply ``learning'' like humans do. However, the sheer scale of this automated extraction essentially commodifies human creativity, transforming it into raw fuel for commercial software products without equitable profit sharing. Furthermore, the ``black box'' nature of neural networks makes it nearly impossible to trace exactly which piece of art contributed to a specific generated output, complicating legal accountability.
%     \item \textbf{The Proliferation of Deepfakes and Synthetic Fraud:} The accessibility of both commercial APIs and uncensored open-source GenAI tools has exponentially lowered the barrier to entry for malicious actors. This has led to an explosion in the creation of ``deepfakes''—synthetic media weaponized to generate non-consensual explicit content (image-based sexual abuse), execute highly convincing financial fraud via voice cloning, and spread targeted political misinformation to manipulate elections. The core issue is that the technology to generate fakes has vastly outpaced the technology to detect them.
% \end{itemize}

% \subsection{Key Stakeholders and Competing Interests}
% Our brainstorming sessions identified four primary stakeholders, each operating under conflicting motivations and pressures:
% \begin{itemize}
%     \item \textbf{AI Developers and Tech Corporations}: This group encompasses data engineers, AI researchers, and corporate executives. Driven by immense pressure from venture capitalists and shareholders, their primary motivation is to win the global ``AI Arms Race.'' They prioritize rapid innovation, model optimization, and maximizing market share, often viewing restrictive data privacy laws as bottlenecks to technological progress.
%     \item \textbf{Original Content Creators and Artists}: Illustrators, photographers, authors, and journalists whose life's work is being ingested by algorithms. They face an existential threat: their own stolen intellectual property is being used to train the very machines that are actively devaluing their labor and threatening their livelihoods.
%     \item \textbf{The Open-Source Community}: Independent developers and researchers who advocate for democratizing AI access. While their goal is to prevent tech monopolies from controlling AI, releasing highly capable, uncensored models to the public inadvertently removes safety guardrails, enabling bad actors to generate harmful content untethered from corporate oversight.
%     \item \textbf{The General Public and Society at Large}: Everyday consumers who benefit from accessible creative tools and increased productivity. However, they are simultaneously the most vulnerable victims of deepfake-driven scams, identity theft, and the psychological toll of a polluted digital ecosystem where reality is constantly questioned.
% \end{itemize}

% \subsection{Constraints and Operational Challenges}
% Engineers attempting to solve this crisis face several severe constraints. The fundamental obstacle is the \textbf{``Pacing Problem''} (often related to the Collingridge Dilemma)—the speed of AI technological advancement vastly outstrips the slow, bureaucratic progression of international copyright legislation and regulatory frameworks. By the time a law is drafted to regulate a specific AI capability, the technology has already evolved past it. Additionally, there is a \textbf{Technical Constraint}: deepfake detection is inherently a reactive, cat-and-mouse game. As soon as a detection algorithm identifies a flaw in AI-generated media, generative models are updated to fix that exact flaw, rendering the detection tool obsolete. Engineers must operate within this chaotic legal and technical gray area, balancing the drive for performance optimization against profound ethical responsibilities.

% \subsection{Long-term Risks and Societal Opportunities}
% \begin{itemize}
%     \item \textbf{Opportunities:} If ethically harnessed, GenAI offers unprecedented opportunities for humanity. It enables rapid prototyping in engineering and design, drastically reduces production costs in the independent entertainment industry, accelerates the creation of customized educational visual aids, and empowers individuals lacking technical artistic skills to express complex ideas visually. It can also revolutionize medical imaging analysis and architectural conceptualization.
%     \item \textbf{Risks:} Beyond the immediate economic risks of widespread copyright lawsuits and job displacement in the creative sector, the ultimate long-term risk is an epistemic crisis—the complete degradation of a shared objective reality. If society reaches a point where it can no longer cryptographically or visually distinguish between an authentic photograph and an AI-generated deepfake, we face the \textbf{``Liar's Dividend.''} This is a psychological phenomenon where the mere existence of deepfakes allows guilty individuals to plausibly deny real, authentic evidence (e.g., video footage of a crime or corruption) by simply claiming it was generated by AI. This threatens to fundamentally compromise journalism, legal evidence, and the foundation of democratic processes.
% \end{itemize}

% \newpage
% \section{Brainstorming \& Problem Framing}

% \subsection{Background and Context of the Problem}
% The rapid expansion of Generative Artificial Intelligence (GenAI) has fundamentally transformed the digital content ecosystem. Tools capable of generating text, images, audio, and video now allow users to produce high-quality content with unprecedented speed, scale, and accessibility. From an engineering and societal perspective, this shift is significant because it lowers creative barriers, increases productivity, and enables new forms of digital innovation. However, these benefits are accompanied by major concerns regarding data provenance, copyright compliance, misuse of synthetic media, and the broader reliability of digital information. \\

% To frame this issue systematically, our team applied structured brainstorming methods, primarily the \textbf{5W1H framework}, supported by collaborative mind mapping. This process helped us move beyond general debate and identify the problem as a socio-technical challenge involving not only technical capability, but also governance, accountability, and public trust. Rather than viewing GenAI solely as a breakthrough or solely as a threat, we approached it as an engineering dilemma in which innovation, legal responsibility, and social impact must be evaluated together.

% \subsection{Structured Brainstorming Summary}
% Using the 5W1H framework, our team identified the following key dimensions of the problem:

% \begin{itemize}
%     \item \textbf{Who:} The issue affects multiple stakeholder groups, including AI developers and technology firms, data engineers, content creators and artists, digital platforms, regulators, and the general public, especially individuals vulnerable to deepfake-related harm.
%     \item \textbf{What:} The central conflict concerns the use of large-scale online data for GenAI training without clear consent, attribution, or compensation, alongside the increasing ability of these systems to generate highly realistic synthetic content that can be misleading or harmful.
%     \item \textbf{When:} The problem is highly urgent in the current phase of GenAI commercialization, where the pace of technological development has exceeded the pace of legal, regulatory, and detection frameworks.
%     \item \textbf{Where:} The issue is most visible across the global digital ecosystem, including social media platforms, online marketplaces, content-sharing environments, and increasingly professional sectors such as education, journalism, and design.
%     \item \textbf{Why:} Organizations seek larger datasets and faster model improvement to maintain technological competitiveness, while creators seek protection of their intellectual labor and society requires trustworthy digital information.
%     \item \textbf{How:} The challenge is not whether GenAI should be used, but how it can be designed, deployed, and governed in ways that preserve innovation while reducing copyright disputes, deception, and misuse.
% \end{itemize}

% \subsection{Core Problem Definition}
% Based on the brainstorming process, the team defines the core problem as follows:

% \begin{tcolorbox}[
%     colback=blue!5!white,
%     colframe=blue!50!black,
%     coltitle=white,        % màu chữ tiêu đề
%     colbacktitle=blue!50!black, % màu nền tiêu đề
%     boxrule=1.5pt,         % độ dày viền
%     enhanced,
%     attach boxed title to top left = {yshift = -2 mm, xshift = 5 mm}, % canh tiêu đề
%     title={\large \textbf{Core problem}}
% ]

% \textit{How can Generative AI systems for digital content creation be developed and governed so that their innovative benefits are preserved while minimizing intellectual property violations, harmful deepfake misuse, and the erosion of trust in digital media?}

% \end{tcolorbox}

% This problem has two closely connected dimensions. First, GenAI models rely on large-scale datasets whose provenance, ownership status, and licensing conditions are often unclear. This creates legal and ethical concerns related to copyright, consent, compensation, and accountability. Second, the same technologies that enable creative productivity can also be used to generate deceptive or malicious synthetic media, including impersonation, non-consensual imagery, fraud, and misinformation. These two dimensions are interconnected because both emerge from the same underlying issue: the rapid deployment of powerful generative systems without sufficiently mature mechanisms for traceability, oversight, and responsible use.

% \subsection{Key Stakeholders and Competing Interests}
% Our brainstorming identified several major stakeholder groups whose interests are partially aligned but often in conflict:

% \begin{itemize}
%     \item \textbf{AI developers and technology companies} seek rapid model improvement, competitive advantage, and commercial scalability. Their priorities often include access to large datasets, reduced friction in deployment, and continued innovation.

%     \item \textbf{Content creators and rights holders}, including artists, writers, photographers, and journalists, seek recognition, consent, fair compensation, and protection against the unauthorized reuse of their work.

%     \item \textbf{Digital platforms and content distributors} are responsible for hosting, distributing, and moderating content. They face pressure to support innovation and user engagement while also reducing abuse, misinformation, and reputational damage.

%     \item \textbf{Regulators, policymakers, and legal institutions} are concerned with enforceability, public safety, copyright protection, and the adaptation of legal frameworks to new technological realities.

%     \item \textbf{The general public} benefits from easier access to creative tools and lower content-production costs, but also faces risks related to deception, identity misuse, privacy invasion, and reduced trust in digital evidence.

%     \item \textbf{Individuals directly targeted by deepfakes} represent a particularly vulnerable group, as they may suffer reputational, psychological, social, or financial harm even when they did not participate in the technology at all.
% \end{itemize}

% These stakeholders illustrate that the issue is not a simple conflict between ``technology'' and ``society.'' Instead, it is a complex engineering governance problem involving trade-offs between openness and safety, innovation and accountability, automation and human rights.

% \subsection{Constraints and Operational Challenges}
% From an engineering and policy perspective, several constraints make this problem difficult to address effectively. \\

% First, there is a \textbf{data provenance constraint}. Once content has been collected at large scale from multiple online sources, it becomes difficult to verify origin, authorship, consent status, and licensing conditions. This limits traceability and complicates compliance. \\

% Second, there is a \textbf{technical constraint} in detecting synthetic media. Deepfake detection remains a reactive process in which detection tools are continuously challenged by improved generation techniques. This creates an ongoing adversarial cycle between generation and detection. \\

% Third, there is a \textbf{regulatory constraint}. Legal frameworks on copyright, privacy, and AI accountability differ across jurisdictions and typically evolve more slowly than the underlying technologies. This creates uncertainty for engineers, organizations, and affected users. \\

% Fourth, there is an \textbf{operational trade-off constraint}. Stronger safeguards---such as watermarking, dataset filtering, or stricter access controls---may improve safety and accountability, but they can also increase cost, reduce usability, slow deployment, or limit openness. As a result, any engineering response must balance performance, compliance, scalability, and risk mitigation.

% \subsection{Risks and Opportunities}
% The brainstorming process also highlighted that GenAI should not be framed only in terms of threats. The technology creates substantial opportunities, but these opportunities can only be sustained if the associated risks are managed responsibly. \\

% \textbf{Opportunities} include faster prototyping, lower barriers to content creation, increased accessibility for users without advanced technical or artistic skills, and productivity gains in sectors such as education, design, communication, and engineering visualization. In this sense, GenAI can expand participation in digital creativity and accelerate innovation. \\

% At the same time, the \textbf{risks} are significant. These include copyright disputes, labor displacement in creative industries, fraudulent impersonation, non-consensual content generation, and large-scale misinformation. At a broader societal level, the most severe risk is the weakening of trust in digital media. If authentic and synthetic content become increasingly difficult to distinguish, public confidence in journalism, legal evidence, and online communication may deteriorate. This creates not only a technical problem, but also a societal and institutional one.

% \subsection{Problem Framing Outcome}
% Overall, our Stage 1 brainstorming indicates that the rise of GenAI in digital content creation should be understood as a \textbf{socio-technical engineering challenge} rather than a purely legal or purely technological issue. The central challenge is to develop systems and governance mechanisms that preserve the value of innovation while ensuring accountability, fairness, and the integrity of digital information. This framing provides the foundation for the later stages of the project, where more specific evaluation and solution directions can be developed.

\newpage
\section{Brainstorming \& Problem Framing}

\subsection{Background and Context of the Problem}
The rapid expansion and commercialization of Generative Artificial Intelligence (GenAI) has fundamentally transformed the digital content ecosystem. Tools capable of generating text, images, audio, and video now allow users to produce high-quality content with unprecedented speed, scale, and accessibility. Examples include image-generation systems such as Midjourney and Stable Diffusion, video-generation models such as Veo and Sora, and increasingly accessible voice-cloning technologies. From both engineering and societal perspectives, this shift is highly significant because it lowers barriers to creativity, increases productivity, and enables new forms of digital innovation. However, these benefits are accompanied by major concerns regarding data provenance, copyright compliance, synthetic media misuse, and the broader reliability of digital information. \\

To frame this issue systematically, our team applied structured brainstorming methods, primarily the \textbf{5W1H framework}, supported by collaborative mind mapping. This process helped us move beyond general discussion and identify the problem as a \textbf{socio-technical engineering challenge} involving not only technical capability, but also governance, accountability, legal uncertainty, and public trust. Rather than viewing GenAI solely as a breakthrough or solely as a threat, we approached it as an engineering dilemma in which innovation, legal responsibility, and social impact must be evaluated together. \\

A key insight from our brainstorming is that the controversy surrounding GenAI does not arise only from its outputs, but also from the way these systems are built and deployed. State-of-the-art models are commonly trained on extremely large datasets collected from online environments at scale, often with limited transparency regarding origin, authorship, consent, or licensing status. This creates a fundamental tension between rapid model improvement and the protection of intellectual labor, while also raising questions about whether existing governance mechanisms are capable of keeping pace with technological change.

\subsection{Structured Brainstorming Summary}
Using the 5W1H framework, our team identified the following key dimensions of the problem:

\begin{itemize}
  \item \textbf{Who:} The issue affects multiple stakeholder groups, including AI developers and technology firms, data engineers, content creators and artists, digital platforms, open-source communities, regulators, and the general public, especially individuals vulnerable to deepfake-related harm.

  \item \textbf{What:} The central conflict concerns the use of large-scale online data for GenAI training without clear consent, attribution, or compensation, alongside the increasing ability of these systems to generate highly realistic synthetic content that may be deceptive, harmful, or socially disruptive.

  \item \textbf{When:} The problem is highly urgent in the current phase of GenAI commercialization, where the pace of technological development has exceeded the pace of legal, regulatory, and detection frameworks.

  \item \textbf{Where:} The issue is most visible across the global digital ecosystem, including social media platforms, online marketplaces, content-sharing environments, and increasingly professional sectors such as education, journalism, design, and communication.

  \item \textbf{Why:} Organizations seek larger datasets and faster model improvement to maintain technological competitiveness, while creators seek protection of their intellectual labor and society requires trustworthy digital information and enforceable accountability.

  \item \textbf{How:} The challenge is not whether GenAI should be used, but how it can be designed, deployed, and governed in ways that preserve innovation while reducing copyright disputes, deception, harmful misuse, and systemic erosion of trust.
\end{itemize}

\subsection{Core Problem Definition}
Based on the brainstorming process, the team defines the core problem as follows:

\begin{tcolorbox}[
    colback=blue!5!white,
    colframe=blue!50!black,
    coltitle=white,
    colbacktitle=blue!50!black,
    boxrule=1.5pt,
    enhanced,
    attach boxed title to top left = {yshift=-2mm, xshift=5mm},
    title={\large \textbf{Core problem}}
  ]
  \textit{How can Generative AI systems for digital content creation be developed and governed so that their innovative benefits are preserved while minimizing intellectual property violations, harmful deepfake misuse, and the erosion of trust in digital media?}
\end{tcolorbox}

This problem has two closely connected dimensions. First, GenAI models rely on large-scale datasets whose provenance, ownership status, and licensing conditions are often unclear. This creates legal and ethical concerns related to copyright, consent, compensation, accountability, and fair value distribution. The issue is made more difficult by the black-box nature of many generative systems, which makes it hard to determine how particular copyrighted works influence specific outputs and therefore complicates traceability and legal responsibility. \\

Second, the same technologies that enable creative productivity can also be used to generate deceptive or malicious synthetic media, including impersonation, non-consensual imagery, voice-cloning scams, fraud, and targeted misinformation. These two dimensions are interconnected because both emerge from the same underlying issue: the rapid deployment of powerful generative systems without sufficiently mature mechanisms for traceability, oversight, and responsible use.

\subsection{Key Stakeholders and Competing Interests}
Our brainstorming identified several major stakeholder groups whose interests are partially aligned but often in conflict:

\begin{itemize}
  \item \textbf{AI developers and technology companies} seek rapid model improvement, competitive advantage, and commercial scalability. Their priorities often include access to large datasets, reduced friction in deployment, and continued innovation under intense market pressure.

  \item \textbf{Content creators and rights holders}, including artists, writers, photographers, and journalists, seek recognition, consent, fair compensation, and protection against the unauthorized reuse of their work in training datasets and generated outputs.

  \item \textbf{Digital platforms and content distributors} are responsible for hosting, distributing, and moderating content. They face pressure to support innovation and user engagement while also reducing abuse, misinformation, fraud, and reputational damage.

  \item \textbf{Open-source AI communities} promote broader access to generative technologies and often resist excessive concentration of AI power in a few corporations. However, the open release of highly capable models may also reduce centralized safety controls and make harmful misuse harder to prevent.

  \item \textbf{Regulators, policymakers, and legal institutions} are concerned with enforceability, public safety, copyright protection, privacy, and the adaptation of legal frameworks to new technological realities.

  \item \textbf{The general public} benefits from easier access to creative tools and lower content-production costs, but also faces risks related to deception, identity misuse, privacy invasion, and declining trust in digital evidence.

  \item \textbf{Individuals directly targeted by deepfakes} represent a particularly vulnerable group, as they may suffer reputational, psychological, social, or financial harm even when they did not participate in the technology at all.
\end{itemize}

These stakeholders illustrate that the issue is not a simple conflict between ``technology'' and ``society.'' Instead, it is a complex engineering governance problem involving trade-offs between openness and safety, innovation and accountability, automation and human rights, and accessibility and control.

\subsection{Constraints and Operational Challenges}
From an engineering and policy perspective, several constraints make this problem difficult to address effectively. \\

First, there is a \textbf{\textit{data provenance constraint}}. Once content has been collected at large scale from multiple online sources, it becomes difficult to verify origin, authorship, consent status, and licensing conditions. This limits traceability and complicates compliance, especially when datasets have already been integrated into large model-training pipelines. \\

Second, there is a \textbf{\textit{technical constraint}} in detecting synthetic media. Deepfake detection remains largely reactive, with detection tools continuously challenged by improved generation techniques. This creates an ongoing adversarial cycle between generation and detection, in which defensive methods often lag behind offensive capabilities. \\

Third, there is a \textbf{\textit{regulatory constraint}}. Legal frameworks on copyright, privacy, and AI accountability differ across jurisdictions and typically evolve more slowly than the underlying technologies. This reflects a broader \textbf{\textit{pacing problem}}, often associated with the \textbf{\textit{Collingridge Dilemma}}, in which technologies become difficult to regulate effectively because their consequences are fully recognized only after widespread deployment. \\

Fourth, there is an \textbf{\textit{operational trade-off constraint}}. Stronger safeguards---such as watermarking, dataset filtering, provenance labeling, identity verification, or stricter access controls---may improve safety and accountability, but they can also increase cost, reduce usability, slow deployment, or limit openness. As a result, any engineering response must balance performance, compliance, scalability, accessibility, and risk mitigation.

\subsection{Risks and Opportunities}
The brainstorming process also highlighted that GenAI should not be framed only in terms of threats. The technology creates substantial opportunities, but these opportunities can only be sustained if the associated risks are managed responsibly. \\

\textbf{Opportunities} include faster prototyping, lower barriers to content creation, increased accessibility for users without advanced technical or artistic skills, and productivity gains in sectors such as education, design, communication, engineering visualization, and creative media production. GenAI can also support personalized educational materials, accelerate concept development, and broaden participation in digital creativity by enabling non-experts to express ideas more effectively. \\

At the same time, the \textbf{risks} are significant. These include copyright disputes, labor displacement in creative industries, fraudulent impersonation, non-consensual synthetic content, large-scale misinformation, and financial scams enabled by realistic voice or image generation. At a broader societal level, one of the most severe risks is the weakening of trust in digital media. If authentic and synthetic content become increasingly difficult to distinguish, public confidence in journalism, legal evidence, and online communication may deteriorate. \\

A particularly serious long-term danger is the \textbf{``liar's dividend''}: the possibility that the existence of deepfakes allows individuals to dismiss genuine evidence as AI-generated. In that scenario, the problem is no longer limited to fake content alone; even authentic recordings, photographs, or testimonies may lose persuasive force. This would create not only a technical challenge, but also an epistemic and institutional one, affecting how truth, evidence, and accountability function in digital society.

\newpage
\subsection{Problem Framing Outcome}
Overall, our Stage 1 brainstorming indicates that the rise of GenAI in digital content creation should be understood as a \textbf{socio-technical engineering challenge} rather than a purely legal or purely technological issue. The central challenge is to develop systems and governance mechanisms that preserve the value of innovation while ensuring accountability, fairness, traceability, and the integrity of digital information. \\

This framing is important because it shifts the discussion away from overly simplistic positions such as ``AI should be fully unrestricted'' or ``AI should be completely blocked.'' Instead, it highlights the need for balanced engineering and governance solutions that can manage trade-offs among innovation, rights protection, public safety, and trust. This problem framing provides a strong foundation for the later stages of the project, where more specific evaluation criteria and solution directions can be developed.
